\section{Исследование влияния полосы обзора (Span)}

В данном разделе исследуется влияние полосы обзора (\textit{Span}) анализатора спектра на визуальное отображение спектрального пика гармонического сигнала при фиксированной центральной частоте настройки $\text{Center Freq} = 1{,}0\text{ ГГц}$ и неизменной полосе пропускания фильтра промежуточной частоты $\text{RBW} = 100\text{ кГц}$.

\subsection{Результаты эксперимента}
Исходное значение полосы обзора, полученное после автоматической настройки, составляло $\text{Span} = 10\text{ МГц}$. Для анализа влияния масштаба по частоте значение полосы обзора было увеличено в 5 раз (до $50\text{ МГц}$), а затем уменьшено в 5 раз относительно исходного (до $2\text{ МГц}$). Значение полосы фильтра $\text{RBW}$ принудительно фиксировалось на исходном уровне 300 кГц. Полученные спектрограммы приведены на рис.~\ref{fig:span_50mhz} и рис.~\ref{fig:span_2mhz}.

\begin{figure}[H]
    \centering
    \begin{minipage}[b]{0.48\textwidth}
        \centering
        % Замените fig4 на имя вашего файла
        \includegraphics[width=\textwidth]{fig4}
        \caption{Полоса обзора увеличена в 5 раз ($\text{Span} = 50\text{ МГц}$)}
        \label{fig:span_50mhz}
    \end{minipage}
    \hfill
    \begin{minipage}[b]{0.48\textwidth}
        \centering
        % Замените fig5 на имя вашего файла
        \includegraphics[width=\textwidth]{fig5}
        \caption{Полоса обзора уменьшена в 5 раз ($\text{Span} = 2\text{ МГц}$)}
        \label{fig:span_2mhz}
    \end{minipage}
\end{figure}

    
    При изменении параметра $\text{Span}$ изменяется масштаб горизонтальной оси экрана (цена одного деления координатной сетки):

    При увеличении $\text{Span}$ до $50\text{ МГц}$ весь экран отображает в 5 раз более широкий диапазон частот, цена деления сетки увеличивается в 5 раз, из-за чего неизменный по частоте отклик в $100\text{ кГц}$ занимает в 5 раз меньше экранных пикселей и визуально кажется узким шпилем;

    При уменьшении $\text{Span}$ до $2\text{ МГц}$ масштаб сетки сильно растягивается (цена деления уменьшается в 5 раз), и тот же самый неизменный частотный отклик занимает значительно большую долю горизонтальной шкалы дисплея, визуально расширяясь.

